% Nouveau EGA V (cote 25): archive cote 25, site page 148 (PDF page 149).
% typescript. Draft compared: none.
% AI-transcribed by Claude Opus 5.5 (claude-opus-5-5, adaptive thinking, effort high), Message Batches API, batch ega25b (msgbatch_01A7n5YuXubUPNK1zsepi8RG), 2026-11-13; original AI transcription, subsequently checked against both supplied scan parts; see PAGES_CHECKED.tsv and CHECK_DECISIONS.tsv.
\scanpage{25/148}
\mspage{\hand{9}}

ment normal de $\over{X}{V}_y$ sur k(y), (puisque de toutes façons \xout{X}  $V_y$ est CM ).
On peut montrer  plus généralement, sous les conditio\unsure{n|g}s de déf.1 , que si
x est un zéro singulier ordinaire, alors x est une singularité isolée
de V($\hand{\phi}$)  en ce sens que toute générisation \xout{de} x' de x dans V($\hand{\phi}$) est un
point régulier de V($\hand{\phi}$)  (donc si le lieu singulier de V($\hand{\phi}$) est fermé, p.
ex. si X est ``ex\xout{\unsure{c}}cellent'' , x est \unsure{bien} un point isolé de ce lieu singuli-%
er). \xout{Lorsque} (Pour bien faire, il faudrait inclure ce résultat au \unsure{Chap|Chaap}
IV , mais  le par. 16 ne semble guère le lieu. Où le mettre ?).  Il
s'ensuit que si X est lisse sur un corps k, et si x est un zéro géométri%
quement singulier ordinaire, alors  x est un point isolé de l'ensemble
\unsure{des|dess} points de non lissité  de V(\over{\unsure{4}}{$\hand{\phi}$}) , donc est \inserted{(géom.)} normal si  $\dim_x(X) \hand{\geqslant} 2$ ,
(mais \inserted[typed]{nécessairement} non normal si $\dim_x X = 1$ , comme on le \deleted{vérifie} \inserted{constate} par exemple sur
le cas type $\hand{\phi} = \unsure{xy}$ , x et y fonctions co\unsure{or|rr}données dans le plan affine).

\hand{Voici le diagramme analogue à (D), mais pour}
\hand{une section $\varphi$ d'un Module inversible K :}

\hand{(D')}
\diagram{
% NOTE: The complete diagram, including all objects, labels and arrows, is handwritten. Horizontal endpoint marks are solid dots; the vertical endpoints are open zeros.
\begin{tikzcd}
 & & \hand{0} & & \\
 & & \hand{\Omega^1_{X/Y}\otimes\underline{O}_{V^1}} \arrow[u, handwritten] & & \\
\hand{\bullet} \arrow[r, handwritten] & \hand{P^1_{X/Y}(K)^{\vee}\otimes\underline{O}_{V^1}} \arrow[r, handwritten] \arrow[ur, "\hand{\mu}", handwritten] & \hand{\Omega^1_{X/S}\otimes\underline{O}_{V^1}} \arrow[r, handwritten] \arrow[u, handwritten] & \hand{\Omega^1_{V^1/S}} \arrow[r, handwritten] & \hand{\bullet} \\
 & \hand{K^{-1}\otimes\underline{O}_{V^1}} \arrow[u, handwritten] \arrow[r, handwritten] & \hand{\Omega^1_{Y/S}\otimes\underline{O}_{V^1}} \arrow[u, handwritten] \arrow[ur, "\hand{\nu}"', handwritten] & & \\
\hand{\omega^{-2}_{X/Y}\otimes K^{\unsure{-2n-1|-n-1}}\otimes\underline{O}_{V^1}} \arrow[ur, handwritten] & & \hand{0} \arrow[u, handwritten] & &
\end{tikzcd}
}

\mspage{\hand{148}}
